%% toc: P3 40 40.7
\chapter{Jump and Complexity Conditions}
\label{sec:p03-w03-p03-070}

For a first-order sentence \(A\) about graphs and an irrational exponent
\(\alpha\in(0,1)\), the number \(f_A(\alpha)\) is the limiting probability that
\(G(n,n^{-\alpha})\) satisfies \(A\)
(\cref{def:p03-w03-p03-069-spectrum}).  The Shelah--Spencer zero--one law
(\cref{thm:p03-w03-p03-065-shelah-spencer}) says that this number is \(0\) or
\(1\).  This chapter asks which \(\{0,1\}\)-valued functions on the
irrationals arise as some \(f_A\).  Spencer records two necessary conditions.
The \emph{Jump Condition} says that, once the quantifier depth \(u\) of \(A\)
is fixed, \(f_A\) can change value only at the points of an explicit
well-ordered set of rationals.  The \emph{Complexity Condition} says that the
left-hand limits of \(f_A\) at rationals \(a/b\), read as a language of unary
words \(0^a1^b\), lie in the polynomial hierarchy
\cite[pp.~18--21]{Spencer2007LogicRandomStructuresAuthor}.

\paragraph{Inputs and conclusions.}
The inputs are a sentence \(A\) of quantifier depth at most \(u\) in the
language of graphs, a real parameter \(\lambda\in(0,1]\), and the theories
\(T^-_\lambda\) of \cref{def:p03-w03-p03-070-theory}.  The conclusions are:
any two models of \(T^-_\lambda\) satisfy the same sentences of depth \(u\)
(\cref{thm:p03-w03-p03-070-transfer}); every sentence of \(T_\alpha\)
holds almost surely in \(G(n,n^{-\alpha})\) for irrational \(\alpha\)
(\cref{thm:p03-w03-p03-070-as-holds}); \(f_A\) is constant between
consecutive points of a splitting set (\cref{thm:p03-w03-p03-070-jump}); the
left limit \(g_A(\lambda)\) exists and equals \(1\) exactly when
\(T^-_\lambda\vdash A\) (\cref{thm:p03-w03-p03-070-left-limit}); each move
of the source's level game is decided by \(T^-_\lambda\)
(\cref{thm:p03-w03-p03-070-dichotomy}); and the
language \(\{0^a1^b:T^-_{a/b}\vdash A\}\) lies in the polynomial hierarchy
(\cref{thm:p03-w03-p03-070-complexity}).

\paragraph{Version followed and what is added.}
The presentation follows the 30-page version of Spencer's chapter
\emph{Logic and Random Structures} hosted on his New York University page,
not the printed chapter in the 2007 Springer volume \emph{Finite Model Theory
and Its Applications}.  Its sections are numbered without a chapter prefix,
so its Sections~10 and~11 are the printed Sections~4.10 and~4.11.  Every
locator below refers to the hosted version.  The source presents both
conditions in outline and says so \cite[p.~21]{Spencer2007LogicRandomStructuresAuthor}.
The definitions, statements and proofs below are written out in full.  Two
points needed repair, and both are stated where they arise.  The
finite-closure bound requires the rounding \(\lfloor x\rfloor+1\) where the
source writes \(\lceil x\rceil\) (\cref{def:p03-w03-p03-070-sequences}).
The source's justification of the Jump Condition relies on a claim that fails
(\cref{thm:p03-w03-p03-070-nonconstant}), so the Jump Condition is proved
here for a refined splitting set \(\mathrm{XPL}^\ast_u\supseteq\mathrm{XPL}_u\).
Whether it holds for the source's \(\mathrm{XPL}_u\) is left open
(\cref{def:p03-w03-p03-070-jump-boundary}).

\section{Weights, closures and the theories \texorpdfstring{\(T^-_\lambda\)}{T-lambda}}
\label{sec:p03-w03-p03-070-closures}

Graphs are irreflexive and symmetric, with adjacency written \(x\sim y\).
Rooted graphs, their parameters \(r,v,e\), and extension statements are as
in \cref{def:p03-w03-p03-065-rooted-graph}.  That definition of dense,
sparse, rigid and safe (\cref{def:p03-w03-p03-065-dense-sparse}) uses an
irrational \(\alpha\).  Rational parameters are needed here, so the notions
are restated for every real \(\lambda\in(0,1]\).  When \(\lambda\) is
irrational the two definitions agree.

At an irrational exponent the \(t\)-closure, the Nonexistence Lemma and the
Finite Closure Lemma are treated in \cref{def:p03-w03-p03-064-closure},
\cref{def:p03-w03-p03-068-closure} and
\cref{thm:p03-w03-p03-068-finite-closure}, and the axioms of \(T_\alpha\)
are listed in \cref{def:p03-w03-p03-068-axioms}.  The statements below
extend these to every real \(\lambda\in(0,1]\).  At a rational \(\lambda\) a
weight can be exactly \(0\), and that case is the subject of this chapter;
at an irrational \(\lambda\) the extended notions coincide with the ones
cited.

\begin{hbdefinition}{def:p03-w03-p03-070-weight}{Weight, density, rigidity and safety at a parameter \(\lambda\)}
Fix a real \(\lambda\in(0,1]\) and a graph \(K\).  Let \(X\subseteq V(K)\)
be arbitrary and let \(Y\subseteq V(K)\setminus X\) be finite and nonempty.
Let \(e_K(X,Y)\) be the cardinality of the set of edges of \(K\) with at
least one end in \(Y\) and both ends in \(X\cup Y\), and put
\[
  w_\lambda(X,Y)\defeq|Y|-\lambda\,e_K(X,Y)
\]
when the edge set is finite.  If it is infinite, put
\(w_\lambda(X,Y)\defeq-\infty\).  Thus weights take values in
\(\mathbb R\cup\{-\infty\}\); sums use the usual extended-real convention,
but no difference of two infinite weights is used.
Then \(Y\) is \emph{\(\lambda\)-dense over \(X\)} if \(w_\lambda(X,Y)<0\), and
\emph{\(\lambda\)-sparse over \(X\)} if \(w_\lambda(X,Y)\geq0\).  Further,
\(Y\) is \emph{\(\lambda\)-rigid over \(X\)} if \(Y\setminus Y_0\) is
\(\lambda\)-dense over \(X\cup Y_0\) for every \(Y_0\subsetneq Y\) (including
\(Y_0=\emptyset\)), and \emph{\(\lambda\)-safe over \(X\)} if every nonempty
\(Y_1\subseteq Y\) is \(\lambda\)-sparse over \(X\).  For a finite set \(F\)
put \(w_\lambda(\emptyset,F)\defeq|F|-\lambda|E(K|_F)|\), with
\(w_\lambda(\emptyset,\emptyset)\defeq0\).  A rooted graph \((R,H)\) is
\(\lambda\)-dense, \(\lambda\)-rigid or \(\lambda\)-safe when
\(V(H)\setminus R\) is so over \(R\) in \(H\); its weight is \(v-e\lambda\).
\end{hbdefinition}

The sets \(X\cup Y_0\) in the definition of rigidity are the source's
nailextensions, and the sets \(X\cup Y_1\) in the definition of safety are
its subextensions \cite[pp.~6--7]{Spencer2007LogicRandomStructuresAuthor}.
For irrational \(\lambda\) no weight vanishes, because \(v\geq1\).  So
\(\lambda\)-sparse means \(v-e\lambda>0\), and the notions coincide with those
of \cref{def:p03-w03-p03-065-dense-sparse}.  For rational \(\lambda\) a weight
can be \(0\), and such an extension is counted as sparse.  This is the
convention the source adopts for its theories \(T^-_\alpha\)
\cite[p.~19]{Spencer2007LogicRandomStructuresAuthor}.

\begin{hblemma}{thm:p03-w03-p03-070-weights}{proved=annex:p03-w03-p03-070-weights}{Calculus of weights}
Let \(\lambda\in(0,1]\) and let \(K\) be a graph.  The base sets \(X,X'\)
below may be arbitrary subsets of \(V(K)\); the extension sets
\(Y,Y_1,Y_2\) are finite and nonempty wherever a weight is taken.
\begin{enumerate}
\item \emph{Additivity.}  If \(X,Y_1,Y_2\) are pairwise disjoint and
      \(Y_1,Y_2\neq\emptyset\), then
      \(w_\lambda(X,Y_1\cup Y_2)=w_\lambda(X,Y_1)+w_\lambda(X\cup Y_1,Y_2)\).
\item \emph{Monotonicity.}  If \(X\subseteq X'\) and \(Y\cap X'=\emptyset\),
      then \(w_\lambda(X',Y)\leq w_\lambda(X,Y)\).
\item \emph{Rigid subextensions.}  If \(Y\) is not \(\lambda\)-safe over
      \(X\), then some nonempty \(Y_1\subseteq Y\) is \(\lambda\)-rigid over
      \(X\).  In particular every \(Y\) that is \(\lambda\)-dense over \(X\)
      contains a nonempty \(\lambda\)-rigid subset over \(X\).
\item \emph{Enlarging the base.}  If \(Y\) is \(\lambda\)-rigid over \(X\) and
      \(X\subseteq X'\), then \(Y\setminus X'\) is empty or
      \(\lambda\)-rigid over \(X'\).
\end{enumerate}
\end{hblemma}

Item~(3) is the source's statement that a nonsafe extension contains a rigid
subextension, and that a dense extension contains a rigid subextension
\cite[pp.~12 and 21]{Spencer2007LogicRandomStructuresAuthor}.  The source
does not prove it; the annex does.  For infinite bases the proof first
reduces each strict weight inequality to finitely many incident edges.
This permits the closure arguments below to apply to arbitrary graphs and
arbitrary intersections of closed sets.

\begin{hbdefinition}{def:p03-w03-p03-070-closure}{Rigid extensions, closed sets and the \(t\)-closure}
Fix \(\lambda\in(0,1]\), a graph \(G\) (finite or infinite) and an integer
\(t\geq0\).  For \(X\subseteq V(G)\), a \emph{rigid \(t\)-extension of
\(X\)} is a set \(Y\subseteq V(G)\setminus X\) with \(1\leq|Y|\leq t\) that
is \(\lambda\)-rigid over \(X\) in \(G\).  The set \(X\) is
\emph{\(t\)-closed} if it has no rigid \(t\)-extension.  The
\emph{\(t\)-closure} \(\operatorname{cl}_t(X)\) is the intersection of all
\(t\)-closed sets containing \(X\).  For a tuple \(\bar x\) of vertices,
\(\operatorname{cl}_t(\bar x)\) is the closure of the set of its entries.
When the graph or the parameter must be named, write
\(\operatorname{cl}^G_t\) or speak of \(\lambda\)-closure.
\end{hbdefinition}

This is the source's second definition of \(\operatorname{cl}_t\), the least
superset with no rigid extension of at most \(t\) vertices.  Its first
definition adds rigid extensions one at a time until none is left
\cite[p.~8]{Spencer2007LogicRandomStructuresAuthor}.  The next lemma shows
that the two agree, a point the source leaves to unstated combinatorial
lemmas.  For \(t=0\) every set is closed and \(\operatorname{cl}_0(X)=X\), as
in the source \cite[p.~9]{Spencer2007LogicRandomStructuresAuthor}.

\begin{hblemma}{thm:p03-w03-p03-070-closure}{proved=annex:p03-w03-p03-070-closure}{Closure is well defined}
Fix \(\lambda\), \(G\) and \(t\), and let \(X\subseteq V(G)\).
\begin{enumerate}
\item \(\operatorname{cl}_t(X)\) is \(t\)-closed, so it is the least
      \(t\)-closed set containing \(X\).
\item If \(X=R_0\subseteq R_1\subseteq\cdots\subseteq R_l\) and each
      \(R_{j+1}\setminus R_j\) is a rigid \(t\)-extension of \(R_j\), then
      \(R_l\subseteq\operatorname{cl}_t(X)\); if moreover \(R_l\) is
      \(t\)-closed, then \(R_l=\operatorname{cl}_t(X)\).
\item If \(t\leq t'\), every \(t'\)-closed set is \(t\)-closed.
\item If \(X\subseteq W\) and \(W\) is \(t\)-closed in \(G\), then
      \(\operatorname{cl}^G_t(X)\subseteq W\) and
      \(\operatorname{cl}^G_t(X)=\operatorname{cl}^{G|_W}_t(X)\).
\item \(X\) has a rigid \(t\)-extension if and only if some \(Y\) with
      \(1\leq|Y|\leq t\) is \(\lambda\)-dense over \(X\).
\end{enumerate}
\end{hblemma}

For \(t\geq1\) put
\[
  \varepsilon_\lambda(t)\defeq\min\Bigl\{\frac{e\lambda-v}{v}:
  1\leq v\leq t,\ e\in\N,\ e\lambda>v\Bigr\}.
\]
For each \(v\) the quotient is least at the least admissible
\(e=\lfloor v/\lambda\rfloor+1\), so the minimum exists and is positive.
It does not increase as \(t\) grows.  This is the source's \(\epsilon\) with
the strict inequality \(v-e\lambda<0\)
\cite[pp.~9 and 20]{Spencer2007LogicRandomStructuresAuthor}.

\begin{hblemma}{thm:p03-w03-p03-070-finite-closure}{proved=annex:p03-w03-p03-070-finite-closure}{Finite Closure Lemma, in counting form}
Let \(\lambda\in(0,1]\), \(t\geq1\), and let \(X\) be a finite set of
vertices of a graph \(G\).  Suppose that no set \(F\subseteq V(G)\) with
\(|F|\leq|X|+\lfloor|X|/\varepsilon_\lambda(t)\rfloor+t\) has
\(w_\lambda(\emptyset,F)<0\).  Then
\[
  |\operatorname{cl}_t(X)\setminus X|\leq\Bigl\lfloor\frac{|X|}{\varepsilon_\lambda(t)}\Bigr\rfloor .
\]
\end{hblemma}

The source states that \(|\operatorname{cl}_t(x_1,\ldots,x_r)|<r+K\) almost
surely, with \(K=\lceil r/\epsilon\rceil\)
\cite[p.~9]{Spencer2007LogicRandomStructuresAuthor}.  For irrational
\(\lambda\) the quotient \(r/\varepsilon_\lambda(t)\) is irrational, and the
two bounds coincide.  For rational \(\lambda\) a chain of extensions of total
weight exactly \(0\) is not excluded, and only the floor bound is available.

\begin{hbdefinition}{def:p03-w03-p03-070-theory}{The theories \(T^-_\lambda\) and \(T_\alpha\)}
Let \(\lambda\in(0,1]\).  For a tuple \(\bar y\) and an integer \(t\geq1\),
\(\bar y\) is \emph{\(t\)-generic over \(\bar x\)} if no set of at most \(t\)
vertices, disjoint from the entries of \(\bar x\bar y\) and
\(\lambda\)-rigid over them, contains a vertex adjacent to an entry of
\(\bar y\).  The theory \(T^-_\lambda\), in the language \(\{\sim\}\) with
equality, consists of the graph axioms and two schemas.
\begin{itemize}
\item \emph{Nonexistence.}  For every finite graph \(F\) with
      \(|V(F)|-\lambda|E(F)|<0\): no copy of \(F\) exists, that is, no
      injective map \(V(F)\to V(G)\) carries edges of \(F\) to edges.
\item \emph{Generic Extension.}  For every \(\lambda\)-safe rooted graph
      \((R,H)\) with roots \(x_1,\ldots,x_r\) and nonroots
      \(y_1,\ldots,y_v\), and every \(t\geq0\): for all distinct
      \(x_1,\ldots,x_r\) there are distinct \(y_1,\ldots,y_v\), distinct from
      the \(x\)'s, such that (1) every edge of \(H\) meeting a \(y\) is
      present, (2) no other pair \(y_iy_j\) or \(x_iy_j\) is an edge, and
      (3) for \(t\geq1\), \(\bar y\) is \(t\)-generic over \(\bar x\).
\end{itemize}
For irrational \(\alpha\in(0,1)\) write \(T_\alpha\defeq T^-_\alpha\).
\end{hbdefinition}

Clause~(3) is first-order: whether a set of \(w\leq t\) further vertices is
\(\lambda\)-rigid over the entries of \(\bar x\bar y\) depends only on the
adjacencies among these \(r+v+w\) vertices, so it is a finite disjunction.
For irrational \(\alpha\), \(T_\alpha\) is the source's almost sure theory
with its two schemas \cite[p.~10]{Spencer2007LogicRandomStructuresAuthor}.
For rational \(\lambda\), \(T^-_\lambda\) is the source's theory obtained by
counting weight-zero rooted graphs as sparse
\cite[p.~19]{Spencer2007LogicRandomStructuresAuthor}.  It states every
property at its natural level of generality, so each property is quantified
over all such rooted graphs.

\begin{hbconvention}{def:p03-w03-p03-070-as}{The almost-sure statement \((\mathrm{AS}_\alpha)\)}
For an irrational \(\alpha\in(0,1)\), \((\mathrm{AS}_\alpha)\) is the
statement that every sentence of \(T_\alpha\) holds almost surely in
\(G(n,n^{-\alpha})\) (\cref{def:p03-w03-p03-065-almost-sure}).  It is the
probabilistic half of the Shelah--Spencer zero--one law, and
\cref{thm:p03-w03-p03-070-as-holds} below proves it for every irrational
\(\alpha\).
\end{hbconvention}

The source proves the Nonexistence schema almost surely by a first-moment
count \cite[p.~8]{Spencer2007LogicRandomStructuresAuthor}.  For Generic
Extension it indicates an argument from a large-deviation estimate that it
does not prove \cite[pp.~10--11]{Spencer2007LogicRandomStructuresAuthor}.
That argument is carried out in full in
\cref{thm:p03-w03-p03-064-theory-as}, for the theory \(T_\alpha\) of
\cref{def:p03-w03-p03-064-theory-alpha}.  That theory is written with rigid
extensions and genericity in the words of the source's Section~5, while
\cref{def:p03-w03-p03-070-theory} is written with weights.  The next lemma
shows that the two presentations have the same axioms, up to equivalence in
every graph.

\begin{hblemma}{thm:p03-w03-p03-070-same-theory}{proved=annex:p03-w03-p03-070-same-theory}{The two presentations of \(T_\alpha\) agree}
Let \(\alpha\in(0,1)\) be irrational.
\begin{enumerate}
\item For finite disjoint vertex sets \(X,Y\) of a graph \(G\) with
      \(Y\neq\emptyset\), \(Y\) is \(\alpha\)-rigid over \(X\)
      (\cref{def:p03-w03-p03-070-weight}) if and only if \(Y\) is rigid over
      \(X\) in the sense of \cref{def:p03-w03-p03-064-rigid-extension}, and
      if and only if an enumeration of \(Y\) forms a rigid extension over an
      enumeration of \(X\) in the sense of that definition.
\item A rooted graph is \(\alpha\)-safe if and only if it is safe in the
      sense of \cref{def:p03-w03-p03-065-dense-sparse}.
\item For distinct vertices \(x_1,\ldots,x_r,y_1,\ldots,y_v\) and
      \(t\geq1\), \(\bar y\) is \(t\)-generic over \(\bar x\) in the sense of
      \cref{def:p03-w03-p03-070-theory} if and only if it is \(t\)-generic
      over \(\bar x\) in the sense of \cref{def:p03-w03-p03-064-generic}.
\item The two theories have the same graph axioms and are indexed by the
      same data: finite graphs \(F\) with \(|V(F)|-\alpha|E(F)|<0\) for
      Nonexistence, and pairs of a safe rooted graph and an integer
      \(t\geq0\) for Generic Extension.  For each datum, a graph satisfies the
      axiom of \cref{def:p03-w03-p03-070-theory} if and only if it satisfies
      the axiom of \cref{def:p03-w03-p03-064-theory-alpha}.  Hence the two
      theories have the same models and the same consequences.
\end{enumerate}
\end{hblemma}

\begin{hbtheorem}{thm:p03-w03-p03-070-as-holds}{proved=annex:p03-w03-p03-070-as-holds}{Every sentence of \(T_\alpha\) holds almost surely}
For every irrational \(\alpha\in(0,1)\), \((\mathrm{AS}_\alpha)\) holds: each
sentence of \(T_\alpha\) (\cref{def:p03-w03-p03-070-theory}) holds almost
surely in \(G(n,n^{-\alpha})\).
\end{hbtheorem}

The same conclusion is reached from the weight form of the axioms in the
proof of \cref{thm:p03-w03-p03-065-shelah-spencer}.  From here on the results
carry only the hypotheses that the source states.  The probabilistic input
enters through \cref{thm:p03-w03-p03-070-as-holds} alone; every other step is
a consequence of the axioms and is proved in full.

\section{Duplicator's closure strategy across two parameters}
\label{sec:p03-w03-p03-070-strategy}

The source proves \(T_\alpha\) complete by a strategy for Duplicator in the
Ehrenfeucht game on two models of \(T_\alpha\).  After each round the
closures of the chosen tuples, at a size bound that shrinks as rounds are
used up, are isomorphic \cite[pp.~11--12]{Spencer2007LogicRandomStructuresAuthor}.
Both conditions of this chapter come from running that strategy between
models of two \emph{different} theories \(T^-_\lambda\) and \(T^-_\mu\).
The strategy uses only the classification of pairs \((v,e)\) with \(v\)
below a bound fixed by \(u\), together with the size bounds of
\cref{thm:p03-w03-p03-070-finite-closure}.

\begin{hbdefinition}{def:p03-w03-p03-070-admissible}{Admissible sequences}
Let \(\lambda,\mu\in(0,1]\) and \(u\geq1\).  Integers
\(0=t_0<t_1=1\leq t_2\leq\cdots\leq t_u\) form an \emph{admissible sequence
for \((\lambda,\mu,u)\)} if
\begin{enumerate}
\item[(A1)] for all integers \(1\leq v\leq t_u\) and \(e\geq0\):
      \(v<e\lambda\) if and only if \(v<e\mu\);
\item[(A2)] for \(1\leq j\leq u-1\) and \(\kappa\in\{\lambda,\mu\}\):
      \(t_{j+1}\geq\lfloor(u-j)/\varepsilon_\kappa(t_j)\rfloor+1\).
\end{enumerate}
\end{hbdefinition}

Condition (A1) makes \(\lambda\)-density and \(\mu\)-density agree for sets
of at most \(t_u\) vertices.  Hence rigidity, safety and \(t\)-closedness for
\(t\leq t_u\) agree as well.

\begin{hblemma}{thm:p03-w03-p03-070-generic-closure}{proved=annex:p03-w03-p03-070-generic-closure}{Closures inside a generic extension}
Let \(\mu\in(0,1]\), \(0\leq t\leq T\), and let \(G'\) be a graph.  Let
\(O\subseteq V(G')\) be finite and \(T\)-closed, and let \(N'\) be a finite
set disjoint from \(O\) such that no set of at most \(t\) vertices outside
\(O\cup N'\) that is \(\mu\)-rigid over \(O\cup N'\) has a vertex adjacent to
\(N'\).  Then
\(\operatorname{cl}^{G'}_t(X)=\operatorname{cl}^{G'|_{O\cup N'}}_t(X)\) for
every \(X\subseteq O\cup N'\).
\end{hblemma}

\begin{hblemma}{thm:p03-w03-p03-070-extension}{proved=annex:p03-w03-p03-070-extension}{One round of Duplicator's strategy}
Let \(t_0,\ldots,t_u\) be admissible for \((\lambda,\mu,u)\), let
\(G\models T^-_\lambda\) and \(G'\models T^-_\mu\), and take closures in
\(G\) at \(\lambda\) and in \(G'\) at \(\mu\).  Let \(1\leq i\leq u\),
\(s\leq u-i\), \(\bar x\in V(G)^s\), \(\bar x'\in V(G')^s\).  Suppose
\(\varphi\) is an isomorphism from \(G|_{\operatorname{cl}_{t_i}(\bar x)}\)
onto \(G'|_{\operatorname{cl}_{t_i}(\bar x')}\) with \(\varphi(x_j)=x'_j\)
for all \(j\).  Then these closures are finite, and for every \(y\in V(G)\)
there are \(y'\in V(G')\) and an isomorphism \(\varphi_1\) from
\(G|_{\operatorname{cl}_{t_{i-1}}(\bar xy)}\) onto
\(G'|_{\operatorname{cl}_{t_{i-1}}(\bar x'y')}\) with \(\varphi_1(x_j)=x'_j\)
and \(\varphi_1(y)=y'\).  The same holds with the roles of
\((G,\lambda)\) and \((G',\mu)\) exchanged.
\end{hblemma}

The proof separates the source's two cases.  If \(y\) lies in
\(\operatorname{cl}_{t_i}(\bar x)\), its partner is its image under
\(\varphi\) (the Inside case).  Otherwise the new part of the smaller
closure is a safe extension of the old closure.  Duplicator copies it in
\(G'\) by the Generic Extension axiom with genericity parameter \(t_{i-1}\)
(the Outside case) \cite[p.~12]{Spencer2007LogicRandomStructuresAuthor}.

\begin{hbtheorem}{thm:p03-w03-p03-070-transfer}{proved=annex:p03-w03-p03-070-transfer}{Transfer between the theories \(T^-_\lambda\) and \(T^-_\mu\)}
Let \(t_0,\ldots,t_u\) be admissible for \((\lambda,\mu,u)\), and let
\(G\models T^-_\lambda\) and \(G'\models T^-_\mu\).  Let \(0\leq i\leq u\),
\(s\leq u-i\), and let \(\bar x,\bar x',\varphi\) be as in
\cref{thm:p03-w03-p03-070-extension}.  Then for every formula
\(\theta(z_1,\ldots,z_s)\) of quantifier depth at most \(i\),
\[
  G\entails\theta[\bar x]\iff G'\entails\theta[\bar x'].
\]
In particular \(G\) and \(G'\) satisfy the same sentences of quantifier depth
at most \(u\).  Consequently, if \(T^-_\lambda\) and \(T^-_\mu\) are both
consistent, they prove the same sentences of quantifier depth at most \(u\).
\end{hbtheorem}

For \(\lambda\in(0,1]\) and \(u\geq1\), the \emph{strict sequence}
\(s^\lambda_0,\ldots,s^\lambda_u\) is defined by \(s^\lambda_0=0\),
\(s^\lambda_1=1\), and
\[
  s^\lambda_{j+1}=\max\bigl(s^\lambda_j,\ \lfloor(u-j)/\varepsilon_\lambda(s^\lambda_j)\rfloor+1\bigr)
  \qquad(1\leq j\leq u-1).
\]
It is admissible for \((\lambda,\lambda,u)\).

\begin{hbcorollary}{thm:p03-w03-p03-070-complete}{proved=annex:p03-w03-p03-070-complete}{Completeness of \(T^-_\lambda\) and the zero--one law}
\begin{enumerate}
\item For every \(\lambda\in(0,1]\), any two models of \(T^-_\lambda\) are
      elementarily equivalent.  Hence \(T^-_\lambda\) is complete if it is
      consistent.
\item Let \(\alpha\in(0,1)\) be irrational.  Then \(T_\alpha\) is
      consistent and complete.
      For every sentence \(A\), \(f_A(\alpha)\) exists, and
      \(f_A(\alpha)=1\) if \(T_\alpha\vdash A\), \(f_A(\alpha)=0\)
      otherwise.
\item For every \(\lambda\in(0,1]\), \(T^-_\lambda\) is consistent and
      complete.  Moreover, every finite subset of \(T^-_\lambda\)
      is contained in \(T_\gamma\) for all irrational \(\gamma\) in some
      interval \((\lambda-\delta,\lambda)\).
\end{enumerate}
\end{hbcorollary}

Item~(2) is the source's derivation of the zero--one law from the
completeness of \(T_\alpha\) \cite[p.~12]{Spencer2007LogicRandomStructuresAuthor}.
For \(T^-_\lambda\) with \(\lambda\) rational, the source asserts
completeness without proof \cite[p.~19]{Spencer2007LogicRandomStructuresAuthor}.

\section{The splitting sets}
\label{sec:p03-w03-p03-070-splitting}

\begin{hbdefinition}{def:p03-w03-p03-070-sequences}{The source's sequence and the set \(\mathrm{XPL}_u\)}
Let \(u\geq1\) and \(\lambda\in(0,1]\).  Put \(t_0=0\) and \(t_1=1\).  For
\(1\leq i\leq u-1\), once \(t_i\) is defined, let
\[
  \varepsilon^{\leq}_\lambda(t_i)\defeq\min\Bigl\{\frac{e\lambda-v}{v}:
  1\leq v\leq t_i,\ e\in\N,\ v-e\lambda\leq0\Bigr\}.
\]
If this minimum is \(0\), the sequence is \emph{not well defined} at
\(\lambda\).  Otherwise
\(t_{i+1}\defeq\max\bigl(t_i,\lceil(u-i)/\varepsilon^{\leq}_\lambda(t_i)\rceil\bigr)\).
Write \(t_i(\lambda)\) when the parameter matters.  Then
\(\mathrm{XPL}_u\) is the set of rationals \(\lambda\in(0,1]\) such that
either the sequence \(t_0,\ldots,t_u\) is not well defined at \(\lambda\),
or it is well defined and the numerator of \(\lambda\) in lowest terms is at
most \(t_u(\lambda)\).

The \emph{refined set} \(\mathrm{XPL}^\ast_u\) is the set of rationals
\(\lambda\in(0,1]\) whose numerator in lowest terms is at most
\(s^\lambda_u\), where \(s^\lambda\) is the strict sequence of
\cref{sec:p03-w03-p03-070-strategy}; that sequence is always defined.
\end{hbdefinition}

The sequence and \(\mathrm{XPL}_u\) are the source's
\cite[pp.~18--19]{Spencer2007LogicRandomStructuresAuthor}.  The source also
writes \(\mathrm{EXP}_u\) for the same set.  For irrational \(\lambda\),
\(\varepsilon^{\leq}_\lambda=\varepsilon_\lambda\) and
\(\lceil x\rceil=\lfloor x\rfloor+1\) for the irrational quotients involved,
so \(t_i(\lambda)=s^\lambda_i\).  For a rational \(\lambda=a/b\) in lowest
terms, an equation \(v=e\lambda\) with \(1\leq v\leq t\) is solvable exactly
when \(a\leq t\).  So the sequence fails to be defined only at rationals whose
numerator is at most some \(t_i\).

\begin{hbtheorem}{thm:p03-w03-p03-070-xpl}{proved=annex:p03-w03-p03-070-xpl}{Structure of the splitting sets}
Let \(u\geq1\) and let \(S\) be \(\mathrm{XPL}_u\) or \(\mathrm{XPL}^\ast_u\).
\begin{enumerate}
\item \(\mathrm{XPL}_u\subseteq\mathrm{XPL}^\ast_u\).
\item \(1\in S\) and \(1/b\in S\) for every integer \(b\geq1\).  Hence \(S\)
      has no least element and \(\inf S=0\).
\item For every \(L\in(0,1]\) there is \(\delta>0\) with
      \(S\cap(L-\delta,L)=\emptyset\).
\item For every \(L\in(0,1)\setminus S\) there is \(\delta>0\) with
      \(S\cap(L,L+\delta)=\emptyset\).
\item \(S\) is well ordered by \(>\): every nonempty subset of \(S\) has a
      greatest element.
\item For \(q\in S\) put \(q^-\defeq\max\bigl(S\cap(0,q)\bigr)\).  The open
      intervals \((q^-,q)\), \(q\in S\), are pairwise disjoint, contain no
      point of \(S\), and cover \((0,1)\setminus S\).
\end{enumerate}
\end{hbtheorem}

Items (5) and (6) are the source's description of \(\mathrm{XPL}_u\): it is
well ordered under \(>\), every element has a largest smaller element
\((a/b)^-\), and the gaps form intervals \(I\)
\cite[p.~19]{Spencer2007LogicRandomStructuresAuthor}.  The source attributes
the well ordering to number theory reminiscent of continued fractions and does
not prove it.  The proof in the annex uses only the local bounds of items (3)
and (4).  The source also mentions an interval from the smallest element of
\(\mathrm{XPL}_u\) down to zero.  By item~(2) there is no smallest element, so
that interval does not occur, and the intervals of item~(6) already cover
\((0,1)\setminus S\).

\begin{hbexample}{def:p03-w03-p03-070-explosion}{A sequence that explodes, and an accumulation point}
Take \(u=5\) and \(\alpha=\frac13+\frac1m\) with \(m\geq9\), so
\(\alpha=(m+3)/(3m)\).  Then \(2\alpha<1<3\alpha\), so for \(v=1\) the least
\(e\) with \(e\alpha\geq1\) is \(3\), with no equality, and
\(\varepsilon^{\leq}_\alpha(1)=3\alpha-1=3/m\).  Hence
\(t_2=\max(1,\lceil4m/3\rceil)=\lceil4m/3\rceil\).  The numerator of
\(\alpha\) is \((m+3)/\gcd(m+3,9)\leq m+3\leq4m/3\), since \(m\geq9\).  So
with \(v\) equal to that numerator and \(e\) equal to the denominator,
\(v\leq t_2\) and \(v-e\alpha=0\).  Hence
\(\varepsilon^{\leq}_\alpha(t_2)=0\), and the sequence is not well defined.
Thus \(\frac13+\frac1m\in\mathrm{XPL}_5\) for every \(m\geq9\).  These points
decrease to \(\frac13\), which belongs to \(\mathrm{XPL}_5\) by
\cref{thm:p03-w03-p03-070-xpl}(2).  So \(\frac13\) is approached by
\(\mathrm{XPL}_5\) from above only, as item~(3) requires.  The source's
instance is \(m=10^6\): \(\varepsilon=3\cdot10^{-6}\),
\(t_2=\lceil\frac43\cdot10^6\rceil=1\,333\,334\), and the numerator
\(10^6+3\) is at most \(t_2\)
\cite[pp.~18--19]{Spencer2007LogicRandomStructuresAuthor}.  The source
gives the numerator as \(10^6+1\).  Since \(10^6\equiv1\pmod3\), the
fraction \((10^6+3)/(3\cdot10^6)\) is already in lowest terms, so the
numerator is \(10^6+3\); the conclusion is unchanged.
\end{hbexample}

\begin{hbtheorem}{thm:p03-w03-p03-070-nonconstant}{proved=annex:p03-w03-p03-070-nonconstant}{The sequences need not be constant between consecutive points of \(\mathrm{XPL}_u\)}
Let \(u=2\).  Then \(2/5\) and \(3/7\) are consecutive points of
\(\mathrm{XPL}_2\), that is \((3/7)^-=2/5\).  For irrational
\(\alpha\in(2/5,3/7)\), \(t_2(\alpha)=5\) when \(\alpha<5/12\) and
\(t_2(\alpha)=4\) when \(\alpha>5/12\).  Moreover
\(5/12\notin\mathrm{XPL}_2\) but \(5/12\in\mathrm{XPL}^\ast_2\), and the sign
of \(5-12\alpha\) changes at \(5/12\) while \(5\leq t_2(\alpha)\) for
\(\alpha<5/12\).
\end{hbtheorem}

\section{The Jump Condition}
\label{sec:p03-w03-p03-070-jump}

The source's argument for the Jump Condition has two parts.  First, inside
each interval \(I\) cut out by \(\mathrm{XPL}_u\), the sequences
\(t_0,\ldots,t_u\) are the same.  Second, a sentence of quantifier depth
\(u\) needs only the notions of safe and dense rooted graphs with at most
\(t_u\) nonroots, and these notions are constant on \(I\)
\cite[p.~19]{Spencer2007LogicRandomStructuresAuthor}.
\Cref{thm:p03-w03-p03-070-nonconstant} shows that the first part fails, and
that the needed notions change inside \(I\).  The repair below makes the
second part precise through admissible sequences.  Constancy is then proved
point by point, and the points where that proof breaks down are exactly
those of \(\mathrm{XPL}^\ast_u\).

\begin{hblemma}{thm:p03-w03-p03-070-local}{proved=annex:p03-w03-p03-070-local}{Local admissibility of the strict sequence}
Let \(\lambda\in(0,1]\), \(u\geq1\), and let \(s=s^\lambda\) be the strict
sequence.
\begin{enumerate}
\item There is \(\delta>0\) such that \(s\) is admissible for
      \((\gamma,\lambda,u)\) for every \(\gamma\in(\lambda-\delta,\lambda)\).
\item If \(\lambda\) is irrational, or rational and not in
      \(\mathrm{XPL}^\ast_u\), there is \(\delta>0\) such that \(s\) is
      admissible for \((\gamma,\lambda,u)\) for every
      \(\gamma\in(\lambda-\delta,\lambda+\delta)\cap(0,1]\).
\end{enumerate}
\end{hblemma}

\begin{hbtheorem}{thm:p03-w03-p03-070-local-constancy}{proved=annex:p03-w03-p03-070-local-constancy}{Local constancy and left limits of \(f_A\)}
Let \(A\) be a sentence of quantifier depth at most \(u\), and let
\(\lambda\in(0,1]\).
\begin{enumerate}
\item There is \(\delta>0\) such that for every irrational
      \(\gamma\in(\lambda-\delta,\lambda)\), \(f_A(\gamma)=1\) if
      \(T^-_\lambda\vdash A\) and \(f_A(\gamma)=0\) otherwise.
\item If \(\lambda\) is irrational, or rational and not in
      \(\mathrm{XPL}^\ast_u\), the same holds for every irrational
      \(\gamma\in(\lambda-\delta,\lambda+\delta)\cap(0,1)\).
\end{enumerate}
\end{hbtheorem}

The source's statement of the Jump Condition reads, in substance: if
\(f=f_A\) for a first-order \(A\), then some \(u\) makes \(f\) constant on
every interval \(I\) defined by the splitting set \(\mathrm{XPL}_u\)
\cite[p.~19]{Spencer2007LogicRandomStructuresAuthor}.  The form proved here
uses the refined set.

\begin{hbtheorem}{thm:p03-w03-p03-070-jump}{proved=annex:p03-w03-p03-070-jump}{Jump Condition for the refined splitting set}
Let \(A\) be a first-order sentence of quantifier depth at most \(u\).  For
every \(q\in\mathrm{XPL}^\ast_u\), \(f_A\) is constant on the irrationals of
\((q^-,q)\), where \(q^-\) is taken in \(\mathrm{XPL}^\ast_u\).  Its value
there is \(1\) if \(T^-_q\vdash A\) and \(0\) otherwise.  Hence, if
\(f=f_A\) for some first-order \(A\), there is a \(u\) such that \(f\) is
constant on each interval defined by the splitting set
\(\mathrm{XPL}^\ast_u\).
\end{hbtheorem}

\begin{hbcorollary}{thm:p03-w03-p03-070-spectrum}{proved=annex:p03-w03-p03-070-spectrum}{Spectra lie in the refined splitting set}
The spectrum \(\operatorname{Sp}(A)\) of a first-order sentence \(A\) of quantifier depth at most
\(u\) (\cref{def:p03-w03-p03-069-spectrum}) is contained in
\(\mathrm{XPL}^\ast_u\).
\end{hbcorollary}

\begin{proofboundary}{def:p03-w03-p03-070-jump-boundary}{What is not established about \(\mathrm{XPL}_u\)}
Two claims of the source are not proved in this chapter.
\begin{enumerate}
\item The Jump Condition for the source's \(\mathrm{XPL}_u\).  It would follow
      from \cref{thm:p03-w03-p03-070-jump} if \(f_A\) had no jump at the
      points of \(\mathrm{XPL}^\ast_u\setminus\mathrm{XPL}_u\), such as
      \(5/12\) for \(u=2\).  At such a point the closure bound fails by one
      on the left side, and the needed density notions change sign
      (\cref{thm:p03-w03-p03-070-nonconstant}).  So the transfer argument of
      \cref{thm:p03-w03-p03-070-transfer} does not apply there.  The source's
      justification rests on the constancy of the sequences, which fails.
      Neither a proof nor a counterexample is given here.
\item A zero--one law at rational exponents.  The source asserts that, for
      rational \(\alpha\) outside the exceptional set, a zero--one law holds
      for sentences of quantifier depth \(u\)
      \cite[p.~19]{Spencer2007LogicRandomStructuresAuthor}.  This is a
      statement about \(G(n,n^{-\alpha})\) at a rational exponent.  It needs
      probabilistic estimates beyond the almost sure validity of
      \(T_\gamma\) at irrational \(\gamma\)
      (\cref{thm:p03-w03-p03-070-as-holds}), and the source gives none.
\end{enumerate}
\end{proofboundary}

\section{Left limits and the theories \texorpdfstring{\(T^-_{a/b}\)}{T-minus a/b}}
\label{sec:p03-w03-p03-070-left-limits}

\begin{hbdefinition}{def:p03-w03-p03-070-g}{The left limit \(g_A\)}
For a sentence \(A\) and \(\lambda\in(0,1]\), let \(g_A(\lambda)\) be the
limit of \(f_A(\lambda-\epsilon)\) as \(\epsilon\) decreases to \(0\) through
values with \(\lambda-\epsilon\) irrational, when this limit exists
\cite[p.~19]{Spencer2007LogicRandomStructuresAuthor}.
\end{hbdefinition}

\begin{hbtheorem}{thm:p03-w03-p03-070-left-limit}{proved=annex:p03-w03-p03-070-left-limit}{The left limit is truth in \(T^-_\lambda\)}
For every sentence \(A\) and every \(\lambda\in(0,1]\), \(T^-_\lambda\) is
consistent and complete, \(g_A(\lambda)\) exists, and
\[
  g_A(\lambda)=1\iff T^-_\lambda\vdash A .
\]
For irrational \(\lambda\), \(g_A(\lambda)=f_A(\lambda)\).  If \(A\) has
quantifier depth at most \(u\) and \(\alpha\) is irrational, then
\(f_A(\alpha)=g_A(q)\), where \(q\) is the least element of
\(\mathrm{XPL}^\ast_u\) above \(\alpha\).
\end{hbtheorem}

The last sentence is the source's remark that \(g_A\) on the rationals
determines \(f_A\) \cite[p.~19]{Spencer2007LogicRandomStructuresAuthor}.  The
least element exists by \cref{thm:p03-w03-p03-070-xpl}(6).

\begin{hbwarning}{def:p03-w03-p03-070-side}{From which side \(T^-_\lambda\) is a limit}
The source calls \(T^-_\alpha\) the limiting theory of \(T_{\alpha+\epsilon}\)
as \(\epsilon\downarrow0\), and says that weight-zero rooted graphs are sparse
there \cite[p.~19]{Spencer2007LogicRandomStructuresAuthor}.  The side is the
other one.  If \(v=e\alpha\) with \(e\geq1\), then
\(v-e(\alpha+\epsilon)=-e\epsilon<0\), which is dense, while
\(v-e(\alpha-\epsilon)=e\epsilon>0\), which is sparse.  Counting weight zero
as sparse matches the exponents just \emph{below} \(\alpha\), and that is the
side on which \(g_A\) takes its limit.  \Cref{thm:p03-w03-p03-070-complete}(3)
makes this precise.  From the right, the limit of \(f_A\) at a rational
\(\lambda\notin\mathrm{XPL}^\ast_u\) is the same by
\cref{thm:p03-w03-p03-070-local-constancy}(2).  At points of
\(\mathrm{XPL}^\ast_u\) it can differ, as the spectra of
\cref{sec:p03-w03-p03-069-spectra} show.
\end{hbwarning}

\section{The Complexity Condition}
\label{sec:p03-w03-p03-070-complexity}

The Complexity Condition concerns the language
\[
  L_A\defeq\{0^a1^b:1\leq a\leq b,\ T^-_{a/b}\vdash A\}.
\]
The source writes it as \(\{0^a1^b:A\in T^-_{a/b}\}\in\mathrm{PH}\).  The
restriction \(a\leq b\) only makes \(a/b\) lie in \((0,1]\); a word of another
shape is not in \(L_A\).  The value \(a/b\) need not be in lowest terms, and
the input length is \(a+b\).  The polynomial hierarchy is as in
\cref{def:p04-polynomial-hierarchy}.  The bound proved below is membership
in the intersection \(\Sigma^p_{u+2}\cap\Pi^p_{u+2}\) of the two quantifier
classes; no claim is made about the oracle class
\(\mathsf P^{\Sigma^p_{u+1}}\).

\begin{hbmethod}{def:p03-w03-p03-070-game}{The source's level game for \(T^-_{a/b}\)}
The source turns the question whether \(T^-_{a/b}\vdash A\) into a game,
for \(A\) written as \(Qx_1\cdots Qx_u\,P\) with \(P\) quantifier-free
\cite[pp.~20--21]{Spencer2007LogicRandomStructuresAuthor}.
\begin{enumerate}
\item \emph{Board.}  Levels \(0,1,\ldots,u\).  Level \(i\) holds the possible
      \(t_i\)-closures of \(x_1,\ldots,x_{u-i}\).  The top level holds only
      \(\emptyset\), which is the initial position.  Level \(0\) holds the
      possible graphs on \(x_1,\ldots,x_u\) with their equalities.
\item \emph{Final positions.}  A level-\(0\) position is winning when \(P\)
      is true there.
\item \emph{Moves.}  A move goes from \(H\) at level \(i\) to \(H'\) at level
      \(i-1\).  It is permitted when \(T^-_{a/b}\) proves that every tuple
      with \(t_i\)-closure \(H\) has an extension \(x_{u-i+1}\) whose
      \(t_{i-1}\)-closure is \(H'\).
\item \emph{Rules.}  In round \(i\), Duplicator moves if \(x_i\) is
      existentially quantified and Spoiler moves if it is universally
      quantified.  Duplicator wins if the final position is winning.
\item \emph{Certificates.}  To show that a move is permitted, Duplicator
      draws \(H\) and \(H'\).  For an Inside move she names the new vertex
      and the new closure.  For an Outside move she gives the retained
      vertices, the new vertices with their edges, and the chain of rigid
      extensions that produces the new closure.  Spoiler may answer with a
      nailing that makes a claimed rigid extension sparse, a dense
      subextension of the new part, or a dense extension inside the picture
      that shows the closure is larger.
\end{enumerate}
\end{hbmethod}

The source infers membership in the hierarchy from three facts.  The
\(t_i\) grow at most like \(b^i\).  Positions have polynomial size.  And the
certificate exchange puts permissibility in the second level
\cite[pp.~20--21]{Spencer2007LogicRandomStructuresAuthor}.  The statements
below make this exact, for arbitrary sentences rather than prenex ones, using
the strict sequence.

\begin{hblemma}{thm:p03-w03-p03-070-sizes}{proved=annex:p03-w03-p03-070-sizes}{Size bounds at \(\lambda=a/b\)}
Let \(1\leq a\leq b\), \(c=a/b\), and \(u\geq1\).  Then
\(\varepsilon_c(t)\geq1/(tb)\) for every \(t\geq1\), and
\(s^c_i\leq(ub+1)^{i-1}\) for \(1\leq i\leq u\).  The numbers
\(s^c_1,\ldots,s^c_u\) can be computed from \((a,b)\) in time polynomial in
\(b\) for fixed \(u\).
\end{hblemma}

\begin{hbdefinition}{def:p03-w03-p03-070-positions}{Positions and permitted moves}
Fix \(c\in(0,1]\), \(u\geq1\), and write \(s_i=s^c_i\).  A \emph{level-\(i\)
position} (\(0\leq i\leq u\)) is a pair \((H,\bar h)\) of a finite graph \(H\)
and a tuple \(\bar h\in V(H)^k\) with \(k\leq u-i\), taken up to isomorphism
fixing the tuple.  For a graph \(G\) and \(\bar x\in V(G)^k\), \((G,\bar x)\)
\emph{realizes} \((H,\bar h)\) at level \(i\) if
\((G|_{\operatorname{cl}_{s_i}(\bar x)},\bar x)\cong(H,\bar h)\), closures
taken at \(c\).  For \(i\geq1\) the move from \((H,\bar h)\) to a
level-\((i-1)\) position \((H',\bar h')\) with \(|\bar h'|=k+1\) is
\emph{permitted} if one of the following holds.
\begin{itemize}
\item \emph{Inside:} for some \(y\in V(H)\),
      \((H',\bar h')\cong(H|_{\operatorname{cl}^H_{s_{i-1}}(\bar hy)},\bar hy)\).
\item \emph{Outside:} there are a finite graph \(K\) containing \(H\) as an
      induced subgraph and \(y\in N\defeq V(K)\setminus V(H)\) such that
      \(1\leq|N|\leq s_i\), \(N\) is \(c\)-safe over \(V(H)\) in \(K\), and
      \((H',\bar h')\cong(K|_{\operatorname{cl}^K_{s_{i-1}}(\bar hy)},\bar hy)\).
\end{itemize}
\end{hbdefinition}

\begin{hblemma}{thm:p03-w03-p03-070-moves}{proved=annex:p03-w03-p03-070-moves}{Permitted moves are exactly the realizable moves}
Let \(G\models T^-_c\), \(1\leq i\leq u\), and let \((G,\bar x)\) realize the
level-\(i\) position \((H,\bar h)\) with \(|\bar x|=k\leq u-i\).  Then
\(|V(H)|\leq k+\lfloor k/\varepsilon_c(s_i)\rfloor\).  For every level-\((i-1)\)
position \((H',\bar h')\) with \(|\bar h'|=k+1\), some \(y\in V(G)\) makes
\((G,\bar xy)\) realize \((H',\bar h')\) if and only if the move from
\((H,\bar h)\) to \((H',\bar h')\) is permitted.
\end{hblemma}

The source defines a move to be permitted when \(T^-_{a/b}\) proves that
every tuple with the old closure has an extension with the new one.  It adds
that completeness could have been shown syntactically, the key fact being
that for any two positions on consecutive levels the theory proves either
this statement or the statement that no such extension exists
\cite[pp.~20--21]{Spencer2007LogicRandomStructuresAuthor}.  To state that
dichotomy, positions must be expressed by formulas.

\begin{hblemma}{thm:p03-w03-p03-070-closure-formula}{proved=annex:p03-w03-p03-070-closure-formula}{Closure types are first-order at every parameter}
Let \(c\in(0,1]\), \(t\geq0\), and let \(H\) be a finite graph with a tuple
\(\bar h\in V(H)^k\), \(k\geq0\).  There is a first-order formula
\(\theta^t_{H,\bar h}(z_1,\ldots,z_k)\) in the language \(\{\sim\}\) with
equality such that, for every graph \(G\) and every \(\bar z\in V(G)^k\),
\[
  G\entails\theta^t_{H,\bar h}[\bar z]\iff
  (G|_{\operatorname{cl}^G_t(\bar z)},\bar z)\cong(H,\bar h),
\]
closures taken at \(c\).  If \(\operatorname{cl}^H_t(\bar h)\neq V(H)\), no
pair \((G,\bar z)\) satisfies the right-hand side, and
\(\theta^t_{H,\bar h}\) is satisfied in no graph.
\end{hblemma}

The formula asserts a copy of \(H\) with \(\bar z\) in the place of \(\bar h\)
and says that no set of at most \(t\) further vertices is dense over the
copy.  At an irrational exponent and for distinct roots the same fact is
\cref{thm:p03-w03-p03-064-closure-fo}.

\begin{hbtheorem}{thm:p03-w03-p03-070-dichotomy}{proved=annex:p03-w03-p03-070-dichotomy}{The syntactic dichotomy for moves}
Let \(c\in(0,1]\), \(u\geq1\), \(s_i=s^c_i\) and \(1\leq i\leq u\).  Let
\(P=(H,\bar h)\) be a level-\(i\) position with \(|\bar h|=k\leq u-i\), and
\(P'=(H',\bar h')\) a level-\((i-1)\) position with \(|\bar h'|=k+1\)
(\cref{def:p03-w03-p03-070-positions}).  Put
\(\theta_P=\theta^{s_i}_{H,\bar h}\) and
\(\theta_{P'}=\theta^{s_{i-1}}_{H',\bar h'}\), closures at \(c\), and
\[
  \sigma^+_{P,P'}\defeq\forall\bar z\,\bigl(\theta_P(\bar z)\to
  \exists y\,\theta_{P'}(\bar z,y)\bigr),\qquad
  \sigma^-_{P,P'}\defeq\forall\bar z\,\bigl(\theta_P(\bar z)\to
  \neg\exists y\,\theta_{P'}(\bar z,y)\bigr).
\]
\begin{enumerate}
\item If the move \(P\to P'\) is permitted, then
      \(T^-_c\vdash\sigma^+_{P,P'}\); otherwise \(T^-_c\vdash\sigma^-_{P,P'}\).
      So \(T^-_c\) proves one of the two sentences.
\item \(T^-_c\) is consistent and complete, and \(P\) is realized in some
      model of \(T^-_c\) if and only if \(T^-_c\vdash\exists\bar z\,\theta_P(\bar z)\).
      For such \(P\), exactly one of \(\sigma^+_{P,P'}\) and
      \(\sigma^-_{P,P'}\) is provable in \(T^-_c\), and
      \(T^-_c\vdash\sigma^+_{P,P'}\) if and only if the move \(P\to P'\) is
      permitted.
\end{enumerate}
\end{hbtheorem}

For \(c=a/b\) and \(k=u-i\), \(\sigma^+_{P,P'}\) is the sentence by which the
source defines a permitted move, and item~(2) is its dichotomy
\cite[pp.~20--21]{Spencer2007LogicRandomStructuresAuthor}.  The source states
the dichotomy for the theories \(T_\alpha\) as well; that case is \(c\)
irrational, where \(T^-_c=T_c\).  Two differences from the source are kept
explicit.  First, the closures use the strict sequence \(s^c\).  At an
irrational \(c\) it is the source's \(t_i\)
(\cref{def:p03-w03-p03-070-sequences}).  For \(T^-_{a/b}\) the source
computes its \(t_i\) from the same strict \(\varepsilon_c\) but rounds
\((u-j)/\varepsilon_c\) up \cite[p.~20]{Spencer2007LogicRandomStructuresAuthor},
so the two sequences agree up to the first index at which that quotient is an
integer.  Second, for a position realized in no model both sentences are
provable, vacuously.  Such positions are never reached: the initial position
\((\emptyset,())\) is realized in every model, and by
\cref{thm:p03-w03-p03-070-moves} a permitted move from a realized position
leads to a realized position.  So on every position that the game reaches,
the source's syntactic definition of a permitted move and
\cref{def:p03-w03-p03-070-positions} agree.

The relation ``permitted'' depends only on \(c\), \(u\), \(i\) and the two
positions, not on a model.  For a formula \(\theta(z_1,\ldots,z_k)\) of
quantifier depth at most \(i\) and a level-\(i\) position \(P=(H,\bar h)\) with
\(|\bar h|=k\), define \(\mathrm{Val}_i(\theta,P)\in\{\text{true},\text{false}\}\)
by recursion on \(\theta\).  For atomic \(\theta\), read the equality or
adjacency of the named entries of \(\bar h\) in \(H\).  Negation and
conjunction act truth-functionally.  For \(\theta=\exists w\,\psi\), where
after renaming \(w\) is \(z_{k+1}\),
\(\mathrm{Val}_i(\theta,P)\) is true if and only if some permitted move
\(P\to P'\) has \(\mathrm{Val}_{i-1}(\psi,P')\) true.  Universal quantifiers
are read as \(\neg\exists\neg\).

\begin{hblemma}{thm:p03-w03-p03-070-truth}{proved=annex:p03-w03-p03-070-truth}{Truth is decided by positions}
Let \(G\models T^-_c\), and let \((G,\bar x)\) realize the level-\(i\) position
\(P\).  For every formula \(\theta(z_1,\ldots,z_k)\) of quantifier depth at
most \(i\), \(G\entails\theta[\bar x]\) if and only if
\(\mathrm{Val}_i(\theta,P)\) is true.  In particular, if \(T^-_c\) is
consistent, a sentence \(A\) of quantifier depth at most \(u\) satisfies
\(T^-_c\vdash A\) if and only if \(\mathrm{Val}_u(A,(\emptyset,()))\) is true.
\end{hblemma}

\begin{hbtheorem}{thm:p03-w03-p03-070-complexity}{proved=annex:p03-w03-p03-070-complexity}{Complexity Condition}
Let \(A\) be a first-order sentence of quantifier depth at most \(u\geq1\).
\begin{enumerate}
\item The relation that holds of \((0^a1^b,i,P,P')\) exactly when \(1\leq i\leq u\)
      and the move \(P\to P'\) is permitted for \(c=a/b\) is in
      \(\Sigma^p_2\), with \(P\) and \(P'\) restricted to positions of size
      polynomial in \(b\).
\item \(L_A\in\Sigma^p_{u+2}\cap\Pi^p_{u+2}\).  In particular
      \(\{0^a1^b:A\in T^-_{a/b}\}\in\mathrm{PH}\).
\end{enumerate}
\end{hbtheorem}

The level \(u+2\) comes from counting the hierarchy levels used by the
recursion for \(\mathrm{Val}\).  The source's count is level \(u\) for the game
and level two for permissibility, without a combined bound
\cite[pp.~20--21]{Spencer2007LogicRandomStructuresAuthor}.

\begin{hbexample}{def:p03-w03-p03-070-two-fifths}{The strict sequence at \(2/5\) for depth three}
Take \(c=2/5\) and \(u=3\).  For \(v\geq1\) the least \(e\) with
\(2e/5>v\) is \(\lfloor5v/2\rfloor+1\).  The quotient
\((2e/5-v)/v\) at that \(e\) equals \(1/(5v)\) for odd \(v\) and \(2/(5v)\)
for even \(v\).  The pair \(v=2,e=5\) has weight exactly \(0\) and is
sparse, so it does not enter.  Hence \(\varepsilon_c(1)=1/5\), and
\(s_2=\max(1,\lfloor2\cdot5\rfloor+1)=11\).  Then
\(\varepsilon_c(11)=1/55\), the bound \(1/(tb)\) of
\cref{thm:p03-w03-p03-070-sizes} with equality, attained at \(v=11\).  So
\(s_3=\max(11,\lfloor1\cdot55\rfloor+1)=56\).  By
\cref{thm:p03-w03-p03-070-moves}, positions at level \(2\) have at most
\(1+55=56\) vertices, and positions at level \(1\) have at most
\(2+\lfloor2\cdot5\rfloor=12\).  Since the numerator \(2\) is at most
\(s_3=56\), \(2/5\in\mathrm{XPL}^\ast_3\).
\end{hbexample}

\begin{proofboundary}{def:p03-w03-p03-070-tardos}{Sufficiency of the two conditions}
The source states that the Jump Condition and the Complexity Condition
characterize the functions \(f_A\).  It says that their necessity has been
shown there in outline, and that sufficiency is technically hard and is due
to G\'abor Tardos \cite[p.~21]{Spencer2007LogicRandomStructuresAuthor}.  No
proof of sufficiency is given or read here.  The source's reference list
includes a paper of Spencer and Tardos as submitted
\cite[p.~30]{Spencer2007LogicRandomStructuresAuthor}.  The necessity proved
in this chapter uses the refined set \(\mathrm{XPL}^\ast_u\) in the Jump
Condition.
\end{proofboundary}

\section{Limitations}
\label{sec:p03-w03-p03-070-limitations}

The results describe \(f_A\) only on irrational exponents.  At a rational
exponent the limit probability is a different object.  It may fail to exist,
and it need not equal \(g_A\).  The Complexity Condition concerns the
theories \(T^-_{a/b}\), that is, left limits, and says nothing directly about
\(G(n,n^{-a/b})\).  The bound \(\Sigma^p_{u+2}\cap\Pi^p_{u+2}\) is for a fixed sentence.
Nothing here gives one hierarchy level for all sentences, and the exponent of
the polynomial grows with \(u\).  The theorems on \(f_A\) and \(g_A\) use
probability only through \cref{thm:p03-w03-p03-070-as-holds}, which rests on
the large-deviation and moment estimates of
\cref{thm:p03-w03-p03-064-large-deviation,thm:p03-w03-p03-064-moment-bound}.
The transfer and complexity arguments themselves are purely about models of
\(T^-_\lambda\) and need no probability.

\section{Relationships}
\label{sec:p03-w03-p03-070-relationships}

The constant-probability zero--one law of \cref{sec:p03-zero-one-law} is the
case in which no exponent varies.  Its almost sure theory is a single
complete theory, so no jump question arises.  Local constancy of \(f_A\) at
an irrational exponent is \cref{thm:p03-w03-p03-068-continuity};
\cref{thm:p03-w03-p03-070-local-constancy} extends it to one-sided
constancy at every real exponent and two-sided constancy off the refined
splitting set.  The sparse regime and the
zero--one law at irrational exponents are introduced in
\cref{thm:p03-w03-p03-065-shelah-spencer}.  The function \(f_A\) and the
spectrum are defined in \cref{def:p03-w03-p03-069-spectrum}, where a sentence
with infinite spectrum is constructed.  By
\cref{thm:p03-w03-p03-070-spectrum} every spectrum of a depth-\(u\) sentence
lies inside the well-ordered set \(\mathrm{XPL}^\ast_u\).  The game behind
both conditions is the closure strategy for the Ehrenfeucht game
(\cref{thm:p03-ef}), here run between models of two different theories.  The
Complexity Condition places a family of complete theories, indexed by
rationals written in unary, inside the polynomial hierarchy of
\cref{def:p04-polynomial-hierarchy}.  The descriptive-complexity
presentations of \cref{sec:p09-presentation-n-c-002} concern definability on
finite structures and do not supply these conditions.

\section{Sources}
\label{sec:p03-w03-p03-070-sources}

Spencer, \emph{Logic and Random Structures}, the 30-page version hosted on his
New York University page: Section~10 (the Jump Condition, pp.~18--19), Section~11
(the Complexity Condition and the attribution to Tardos, pp.~19--21), and the
earlier material on closure and the almost sure theory (pp.~6--12) that both
sections rely on \cite{Spencer2007LogicRandomStructuresAuthor}.  The proofs in
the annex are written for this chapter.  Where they depart from the source's
outline, the departure is stated at the corresponding result.
